\chapter{Virtual Memory}\label{ch:b3:virtual-memory}

Every process on a laptop believes it owns the whole memory, from address 0 upwards, and two processes can both store at address \code{0x10000000} without disturbing each other. Neither is true of the hardware: there is one physical memory, shared by all of them and by the kernel. This chapter builds the translation that makes the belief hold, the \omterm{def:b3:virtual-memory:paging}{page tables}, walks it in C exactly as a RISC-V processor walks it, measures the cache that makes it fast, and then uses it for more than isolation: running programs larger than the memory, sharing \omterm{def:b3:virtual-memory:paging}{pages} between processes until one of them writes, and reading files with ordinary loads and stores. Along the way it meets one of the strangest facts of systems, a policy that faults more when given more memory, and proves which policies cannot.

\begin{recall}
In the Year~2 volume a cache keeps copies of recently used blocks of memory close to the processor, and is fast because programs have locality: they reuse what they used recently, and use what lies next to it. When a cache is full, a replacement policy chooses the block to evict; the least recently used (LRU) policy evicts the block unused for longest. A process is created by \code{fork}, which makes a copy of the calling process.
\end{recall}

\section{Address spaces}\label{sec:b3:virtual-memory:spaces}

\begin{definition}[Virtual memory, address space, virtual and physical addresses]\label{def:b3:virtual-memory:space}
\emph{Virtual memory}\index{virtual memory} gives each process its own \emph{address space}\index{address space}, the set of addresses it can use, and translates each address it issues, a \emph{virtual address}\index{virtual address}, into a \emph{physical address}\index{physical address}, a location in the memory chips, by a mapping the kernel controls and the processor applies on every access.
\end{definition}

The translation does three jobs at once. It isolates: a process can only reach the physical memory its mapping includes, so it cannot read or damage another's. It relocates: every program can be linked to run at the same addresses, wherever its memory really is. And it lets the kernel lie usefully: an address may be mapped to nothing yet, and get memory only when it is first used, or be mapped to the same physical memory as an address of another process.

\begin{definition}[Page, page frame, page table, page-table entry, multi-level page table]\label{def:b3:virtual-memory:paging}
\emph{Paging} divides the \omterm{def:b3:virtual-memory:space}{virtual address} space into \emph{pages}\index{page}, blocks of a fixed size, and the physical memory into \emph{page frames}\index{page frame} of the same size. A \emph{page table}\index{page table} maps each virtual page number to a frame number; its entries, \emph{page-table entries}\index{page-table entry}, also carry permission bits (readable, writable, executable, accessible from user mode) and a valid bit. A \emph{multi-level page table}\index{multi-level page table} is a tree of small tables: the high bits of the page number index the root, whose entry points to a table indexed by the next bits, and so on.
\end{definition}

The offset within a \omterm{def:b3:virtual-memory:paging}{page} is not translated: a \omterm{def:b3:virtual-memory:space}{virtual address} is a \omterm{def:b3:virtual-memory:paging}{page} number and an offset, and its \omterm{def:b3:virtual-memory:space}{physical address} is the frame number and the same offset. An older scheme, \emph{segmentation}, translated variable-sized segments (code, data, stack) by adding a base and checking a limit; it needs no tables, but the segments must each find a contiguous hole in physical memory, with the \omterm{def:b3:garbage-collection:fragmentation}{external fragmentation} of chapter~\ref{ch:b3:garbage-collection}. \omterm{def:b3:virtual-memory:paging}{Pages} of one size never fragment physical memory that way: any free frame fits any \omterm{def:b3:virtual-memory:paging}{page}.

\begin{omfigure}
\begin{tikzpicture}[font=\scriptsize, x=10mm, y=5mm, >=Stealth,
  pg/.style={draw, minimum width=14mm, minimum height=5mm, inner sep=0pt}]
  \node[font=\small] at (0.7, 6.2) {process A};
  \node[font=\small] at (4.2, 6.2) {physical frames};
  \node[font=\small] at (7.7, 6.2) {process B};
  \foreach \y/\l/\c in {5/code/omDef!12, 4/data/omThm!12, 3/{}/white, 2/{}/white, 1/stack/omMeth!12}
    \node[pg, fill=\c] (a\y) at (0.7, \y) {\l};
  \foreach \y/\l/\c in {5/code/omDef!12, 4/{}/white, 3/data/omThm!25, 2/{}/white, 1/stack/omMeth!25}
    \node[pg, fill=\c] (b\y) at (7.7, \y) {\l};
  \foreach \y/\l/\c in {6/{}/white, 5/{}/white, 4/A data/omThm!12, 3/code/omDef!12, 2/B stack/omMeth!25, 1/B data/omThm!25, 0/A stack/omMeth!12}
    \node[pg, fill=\c] (f\y) at (4.2, \y - 0.5) {\l};
  \draw[->] (a5.east) -- (f3.west);
  \draw[->] (a4.east) -- (f4.west);
  \draw[->] (a1.east) -- (f0.west);
  \draw[->] (b5.west) -- (f3.east);
  \draw[->] (b3.west) -- (f1.east);
  \draw[->] (b1.west) -- (f2.east);
  \node[align=center, font=\scriptsize, black!60] at (0.7, -0.4) {unmapped pages:\\no frame};
\end{tikzpicture}
\par\smallskip{\small Two \omterm{def:b3:virtual-memory:space}{address spaces} over one physical memory. Each \omterm{def:b3:virtual-memory:paging}{page} of a process maps to any free frame; the two copies of the same program share the frame of its code, which neither can write; \omterm{def:b3:virtual-memory:paging}{pages} that are not mapped use no memory at all.}
\end{omfigure}

\section{Paging and multi-level page tables}\label{sec:b3:virtual-memory:tables}

RISC-V's Sv32 mode, for 32-bit processors, uses \omterm{def:b3:virtual-memory:paging}{pages} of 4~KiB. A \omterm{def:b3:virtual-memory:space}{virtual address} is a 20-bit virtual \omterm{def:b3:virtual-memory:paging}{page} number, split into two 10-bit indices, and a 12-bit offset. % ledger: b3:sv32-va
The 20-bit \omterm{def:b3:virtual-memory:paging}{page} number is translated into a 22-bit physical \omterm{def:b3:virtual-memory:paging}{page} number, so a 32-bit processor can address more than 4~GiB of physical memory. % ledger: b3:sv32-ppn
Each \omterm{def:b3:virtual-memory:paging}{page table} holds $2^{10}$ entries of four bytes, exactly one \omterm{def:b3:virtual-memory:paging}{page}, and the physical \omterm{def:b3:virtual-memory:paging}{page} number of the root table is held in the register \code{satp}. % ledger: b3:sv32-table

\begin{omfigure}
\begin{tikzpicture}[font=\scriptsize, x=3.2mm, y=6mm, >=Stealth]
  \foreach \a/\b/\n/\l/\c in {0/10/{VPN[1]}/{31--22}/omDef!15, 10/20/{VPN[0]}/{21--12}/omThm!15, 20/32/{offset}/{11--0}/black!6}
    { \draw[fill=\c] (\a, 0) rectangle (\b, 1);
      \node at ({(\a + \b) / 2}, 0.5) {\n};
      \node[above, black!60] at ({(\a + \b) / 2}, 1) {bits \l}; }
  \node[anchor=east] at (-0.3, 0.5) {virtual address};
  \node[draw, fill=black!6, minimum width=14mm] (satp) at (-4, -2.2) {\code{satp}};
  \draw[fill=omDef!8] (2, -4.6) rectangle (8, -1.6); \node[below] at (5, -4.6) {root table};
  \draw[fill=omDef!25] (2, -3.4) rectangle (8, -2.9); \node[font=\tiny] at (5, -3.15) {pointer};
  \draw[fill=omThm!8] (14, -4.6) rectangle (20, -1.6); \node[below] at (17, -4.6) {level-0 table};
  \draw[fill=omThm!30] (14, -2.6) rectangle (20, -2.1); \node[font=\tiny] at (17, -2.35) {leaf};
  \draw[fill=black!5] (26, -4.6) rectangle (32, -1.6); \node[below] at (29, -4.6) {page frame};
  \draw[fill=black!20] (26, -3.9) rectangle (32, -3.5); \node[font=\tiny] at (29, -3.7) {the byte};
  \draw[->] (satp.east) -- (2, -2.2);
  \draw[->, omDef] (5, 0) -- (5, -2.9);
  \draw[->, omThm] (15, 0) -- (15, -1.6);
  \draw[->, black!60] (26, 0) -- (27, -3.5);
  \draw[->] (8, -3.15) -- (14, -3.15);
  \draw[->] (20, -2.35) -- (26, -2.35);
\end{tikzpicture}
\par\smallskip{\small An Sv32 translation. \code{satp} locates the root table; VPN[1] selects one of its 1,024 entries, which points to a level-0 table; VPN[0] selects the entry there, a leaf, which holds the frame and the permissions; the offset selects the byte in the frame. Two memory reads before every access the program makes.}
\end{omfigure}

An entry is valid when its bit V is set. If its bits R, W and X are all 0, it points to the next level; otherwise it is a leaf, and they say whether the \omterm{def:b3:virtual-memory:paging}{page} may be read, written and executed. A \omterm{def:b3:virtual-memory:paging}{page} cannot be writable without being readable. % ledger: b3:sv32-rwx
The entry's physical \omterm{def:b3:virtual-memory:paging}{page} number fills its top 22 bits; below it come two bits left to software, then D, A, G, U, X, W, R and V. % ledger: b3:sv32-pte
U allows access from user mode, and G marks a mapping present in every \omterm{def:b3:virtual-memory:space}{address space}. A, the \emph{accessed} bit, records that the \omterm{def:b3:virtual-memory:paging}{page} has been used since A was last cleared, and D, the \emph{dirty} bit, that it has been written. % ledger: b3:sv32-ad
A leaf may also appear at the root, mapping a whole 4~MiB \emph{megapage} at once. % ledger: b3:sv32-megapage

\Cref{lst:b3:virtual-memory:map,lst:b3:virtual-memory:walk} implement the tables over a byte array of 1~MiB that plays the physical memory: \code{map} installs one \omterm{def:b3:virtual-memory:paging}{page}, creating its level-0 table when the first \omterm{def:b3:virtual-memory:paging}{page} of a 4~MiB region is mapped; \code{translate} walks the tables as the specification describes, printing each step.

\omcode{code/bachelor-3/11-virtual-memory/listings/sv32.c}{16}{29}{The indices and \code{map}. An entry is a frame number shifted left by 10, or-ed with its bits.}\label{lst:b3:virtual-memory:map}

\omcode{code/bachelor-3/11-virtual-memory/listings/sv32.c}{31}{70}{The walk. A missing entry or a forbidden access is a \omterm{def:b3:virtual-memory:fault}{page fault}; a successful access sets A, and D for a store.}\label{lst:b3:virtual-memory:walk}

\omcode{code/bachelor-3/11-virtual-memory/listings/sv32_demo.c}{7}{22}{A process of three \omterm{def:b3:virtual-memory:paging}{pages}, and four accesses.}\label{lst:b3:virtual-memory:demo}

\omoutput{code/bachelor-3/11-virtual-memory/outputs/sv32_demo.txt}{Three translations and two faults. The store into the code \omterm{def:b3:virtual-memory:paging}{page} fails on the permissions, after a walk that succeeded; the load at \code{0x40000000} stops at the root. The fetch set the A bit: the code \omterm{def:b3:virtual-memory:paging}{page}'s entry has become \code{0x85b}.}\label{out:b3:virtual-memory:demo}

Why two levels? A one-level table would need an entry for each of the $2^{20}$ \omterm{def:b3:virtual-memory:paging}{pages}, 4~MiB per process, mapped or not. With two, a region of the \omterm{def:b3:virtual-memory:space}{address space} that is not used costs one empty entry in the root, and level-0 tables exist only for the 4~MiB regions that hold something. \Cref{out:b3:virtual-memory:size} counts the cost for three layouts.

\omcode{code/bachelor-3/11-virtual-memory/listings/table_size.py}{1}{29}{Page-table memory for three layouts: a small program, a large array, and \omterm{def:b3:virtual-memory:paging}{pages} scattered one per 4~MiB.}\label{lst:b3:virtual-memory:size}

\omoutput{code/bachelor-3/11-virtual-memory/outputs/table_size.txt}{Tables, counting the root, and their bytes, in total and per mapped \omterm{def:b3:virtual-memory:paging}{page}.}\label{out:b3:virtual-memory:size}

The small program, with code, heap, libraries and stack far apart, needs five tables, 20~KiB, against the 4~MiB of one level. A large dense array costs 4 bytes per \omterm{def:b3:virtual-memory:paging}{page}, the size of its entries. The scattered layout is the worst case: each \omterm{def:b3:virtual-memory:paging}{page} brings a table of its own, 4,160 bytes of tables per \omterm{def:b3:virtual-memory:paging}{page} of data. With 64-bit addresses the gap grows: RISC-V's Sv39 uses three levels for a 39-bit \omterm{def:b3:virtual-memory:space}{address space}, % ledger: b3:sv39
and its Sv48 and Sv57 modes four and five. % ledger: b3:sv48-sv57
 Some processors instead used an \emph{inverted} \omterm{def:b3:virtual-memory:paging}{page table}, one entry per physical frame, searched by hashing the virtual \omterm{def:b3:virtual-memory:paging}{page} number: its size follows the physical memory, not the \omterm{def:b3:virtual-memory:space}{address space}, but sharing a frame between two processes becomes awkward.

\section{The translation cache}\label{sec:b3:virtual-memory:tlb}

Two extra memory reads on every access would make a program that reads memory at every step up to three times slower. Programs have locality, and the same few \omterm{def:b3:virtual-memory:paging}{pages} are translated again and again; the processor keeps the recent translations in a small cache.

\begin{definition}[Translation lookaside buffer, TLB reach, address-space identifier]\label{def:b3:virtual-memory:tlb}
A \emph{translation lookaside buffer}\index{translation lookaside buffer} (TLB\index{TLB}) is a cache of recent translations, from virtual \omterm{def:b3:virtual-memory:paging}{page} number to \omterm{def:b3:virtual-memory:paging}{page-table entry}, searched on every access. On a hit the access uses the cached translation; on a miss the processor (or, on some processors, the kernel) walks the \omterm{def:b3:virtual-memory:paging}{page table} and stores the result. The \emph{TLB reach}\index{TLB reach} is the memory its entries cover at once: the number of entries times the \omterm{def:b3:virtual-memory:paging}{page} size. An \emph{address-space identifier}\index{address-space identifier} (ASID) is a number tagging each entry with the \omterm{def:b3:virtual-memory:space}{address space} it belongs to.
\end{definition}

With a \omterm{def:b3:virtual-memory:tlb}{TLB} hit rate $h$, memory access time $m$ and a walk of $L$ reads on a miss, the average time of an access is
\[ m + (1 - h)\, L\, m, \]
so with two levels and $h = 0.99$ an access costs 2~\% more than without translation. The hit rate depends on the order of the accesses. \Cref{lst:b3:virtual-memory:tlb} simulates a \omterm{def:b3:virtual-memory:tlb}{TLB} of 16 to 128 entries while a $1{,}024 \times 1{,}024$ matrix of 4-byte numbers is summed, each row filling exactly one \omterm{def:b3:virtual-memory:paging}{page}.

\omcode{code/bachelor-3/11-virtual-memory/listings/tlb.c}{10}{25}{A fully associative \omterm{def:b3:virtual-memory:tlb}{TLB} with LRU replacement: \code{access_page} hits or refills one entry.}\label{lst:b3:virtual-memory:tlb}

\omoutput{code/bachelor-3/11-virtual-memory/outputs/tlb.txt}{\omterm{def:b3:virtual-memory:tlb}{TLB} hit rates for three orders of the same sum, by number of entries.}\label{out:b3:virtual-memory:tlb}

By rows, each \omterm{def:b3:virtual-memory:paging}{page} is used 1,024 times in a row: one miss per \omterm{def:b3:virtual-memory:paging}{page}, whatever the size. By columns, each access touches the next row, a new \omterm{def:b3:virtual-memory:paging}{page}, and the 1,024 \omterm{def:b3:virtual-memory:paging}{pages} come back only after all the others: every access misses, even with 128 entries, because the \omterm{def:b3:virtual-memory:paging}{pages} in use far exceed the reach, 512~KiB. In strips of 32 rows, the program reuses 32 \omterm{def:b3:virtual-memory:paging}{pages} before moving on: a \omterm{def:b3:virtual-memory:tlb}{TLB} of 32 entries or more holds them all, and one of 16 misses every time. Large \omterm{def:b3:virtual-memory:paging}{pages} raise the reach: one megapage entry covers as much as 1,024 small ones.

A \omterm{def:b3:kernel-and-scheduling:switch}{context switch} changes \code{satp}, and with it every translation. Without tags the kernel must flush the whole \omterm{def:b3:virtual-memory:tlb}{TLB}, and the next process starts with misses. With an ASID, kept in \code{satp} beside the root table and up to 9 bits wide in Sv32, entries of several processes coexist. % ledger: b3:rv-asid
When the kernel changes a \omterm{def:b3:virtual-memory:paging}{page table}, it must also tell the processor: on RISC-V the instruction \code{sfence.vma} invalidates cached translations, for one address or all, one ASID or all. % ledger: b3:rv-sfence

\section{Demand paging and page replacement}\label{sec:b3:virtual-memory:replacement}

\begin{definition}[Page fault, demand paging]\label{def:b3:virtual-memory:fault}
A \emph{page fault}\index{page fault} is the exception the processor raises when a translation fails: no valid entry, or an access the entry does not allow. With \emph{demand paging}\index{demand paging}, the kernel maps a \omterm{def:b3:virtual-memory:paging}{page} only when the program first touches it: the fault is the request.
\end{definition}

The kernel's fault handler takes the faulting address from a register the trap filled, \code{stval} on RISC-V, % ledger: b3:rv-stval
finds which region of the process it belongs to, and decides. An address in no region is an error: the process receives a signal, usually fatal, the familiar segmentation fault. Otherwise the handler finds a free frame, fills it (with zeros for a new heap or stack \omterm{def:b3:virtual-memory:paging}{page}, from the file for a \omterm{def:b3:virtual-memory:paging}{page} of code or of a mapped file, from the swap area for a \omterm{def:b3:virtual-memory:paging}{page} evicted earlier), writes the entry, and returns to the faulting instruction, which runs again and succeeds. A fault that needs no disk is a \emph{minor} fault; one that must wait for the disk is a \emph{major} fault, and costs milliseconds instead of microseconds.

When no frame is free, one must be taken from some \omterm{def:b3:virtual-memory:paging}{page}, written back to disk first if it is dirty. Which \omterm{def:b3:virtual-memory:paging}{page}? The choice is a replacement policy, as for a cache, but \omterm{def:b3:virtual-memory:fault}{page faults} cost so much that the kernel can afford to think. LRU, recalled above, applies to \omterm{def:b3:virtual-memory:paging}{pages} unchanged.

\begin{definition}[Optimal replacement, clock algorithm, second-chance bit]\label{def:b3:virtual-memory:replacement}
\emph{Optimal replacement}\index{optimal replacement} (OPT) evicts the \omterm{def:b3:virtual-memory:paging}{page} whose next use lies furthest in the future. The \emph{clock algorithm}\index{clock algorithm} keeps the frames in a circle with a hand; each frame has a \emph{second-chance bit}\index{second-chance bit}, set when its \omterm{def:b3:virtual-memory:paging}{page} is used. To evict, the hand advances: a frame with the bit set has it cleared and is passed over; the first frame with the bit clear is evicted.
\end{definition}

OPT needs the future, so no kernel can run it; it is the yardstick that others are measured against, and it is optimal (exercise~\ref{exo:b3:virtual-memory:7}). LRU needs a timestamp on every access, which only the hardware could keep. Clock needs one bit per \omterm{def:b3:virtual-memory:paging}{page}, which the hardware sets for free: it is the A bit of the \omterm{def:b3:virtual-memory:paging}{page-table entry}. It approximates LRU: a \omterm{def:b3:virtual-memory:paging}{page} used since the hand last passed survives one more turn.

\omcode{code/bachelor-3/11-virtual-memory/listings/replace.py}{5}{27}{FIFO evicts the \omterm{def:b3:virtual-memory:paging}{page} loaded first; LRU the \omterm{def:b3:virtual-memory:paging}{page} used least recently.}\label{lst:b3:virtual-memory:fifo}

\omcode{code/bachelor-3/11-virtual-memory/listings/replace.py}{29}{54}{Clock and OPT.}\label{lst:b3:virtual-memory:clock}

\omoutput{code/bachelor-3/11-virtual-memory/outputs/replace.txt}{Faults of each policy with 1 to 5 frames, on two reference strings.}\label{out:b3:virtual-memory:replace}

The first string is the one L.~A. Belady, R.~A. Nelson and G.~S. Shedler published in 1969. % ledger: b3:belady-nelson-shedler-1969
FIFO faults 9 times with 3 frames and 10 times with 4. More memory, more faults.

\begin{definition}[Belady's anomaly, stack algorithm]\label{def:b3:virtual-memory:belady}
\emph{Belady's anomaly}\index{Belady's anomaly} is a reference string on which a policy makes more faults with $f + 1$ frames than with $f$. A policy is a \emph{stack algorithm}\index{stack algorithm} if, on every reference string and at every moment, the set of \omterm{def:b3:virtual-memory:paging}{pages} it holds with $f$ frames is contained in the set it holds with $f + 1$ frames.
\end{definition}

\begin{omfigure}
\begin{tikzpicture}
  \begin{axis}[omrecord, width=0.72\linewidth, height=5.2cm, xmin=0.7, xmax=5.3, ymin=4, ymax=13,
      xtick={1, 2, 3, 4, 5}, xlabel={frames}, ylabel={page faults}, legend pos=south west,
      legend cell align=left]
    \addplot[omThm, thick, mark=*] table[x=frames, y=FIFO]
      {code/bachelor-3/11-virtual-memory/figures/faults.dat};
    \addlegendentry{FIFO (and clock)}
    \addplot[omDef, thick, mark=square*] table[x=frames, y=LRU]
      {code/bachelor-3/11-virtual-memory/figures/faults.dat};
    \addlegendentry{LRU}
    \addplot[omMeth, thick, mark=triangle*] table[x=frames, y=OPT]
      {code/bachelor-3/11-virtual-memory/figures/faults.dat};
    \addlegendentry{OPT}
    \node[omThm, font=\scriptsize, anchor=south] at (axis cs:4, 10.2) {anomaly: 9 $\to$ 10};
  \end{axis}
\end{tikzpicture}
\par\smallskip{\small Faults on Belady's string 1 2 3 4 1 2 5 1 2 3 4 5, from \cref{out:b3:virtual-memory:replace}. LRU and OPT never fault more with more frames; FIFO, and clock with it, fault once more with 4 frames than with 3.}
\end{omfigure}

\begin{theorem}[Stack algorithms have no anomaly]\label{thm:b3:virtual-memory:stack}
A \omterm{def:b3:virtual-memory:belady}{stack algorithm} never makes more faults with $f + 1$ frames than with $f$, on any reference string.
\end{theorem}

\begin{proof}
A reference that hits with $f$ frames finds its \omterm{def:b3:virtual-memory:paging}{page} among the \omterm{def:b3:virtual-memory:paging}{pages} held with $f$ frames, hence, by inclusion, among those held with $f + 1$: it hits there too. So the faults with $f + 1$ frames are among the faults with $f$.
\end{proof}

LRU is a \omterm{def:b3:virtual-memory:belady}{stack algorithm}: with $f$ frames it holds the $f$ \omterm{def:b3:virtual-memory:paging}{pages} most recently used, which are among the $f + 1$ most recently used (exercise~\ref{exo:b3:virtual-memory:10} writes the proof in full). OPT is one too. FIFO is not: after 1~2~3~4~1~2~5, FIFO holds $\{1, 2, 5\}$ with 3 frames and $\{2, 3, 4, 5\}$ with 4. \omterm{def:b3:virtual-memory:belady}{Stack algorithms} were defined by R.~L. Mattson, J. Gecsei, D.~R. Slutz and I.~L. Traiger in 1970, in a paper on evaluating storage hierarchies. % ledger: b3:mattson-1970

\begin{definition}[Working set, thrashing]\label{def:b3:virtual-memory:working-set}
The \emph{working set}\index{working set} of a process at time $t$, for a window $\Delta$, is the set of \omterm{def:b3:virtual-memory:paging}{pages} it used during the last $\Delta$ references. \emph{Thrashing}\index{thrashing} is the collapse of a system whose processes do not all have their working sets in memory: each evicts \omterm{def:b3:virtual-memory:paging}{pages} another will need at once, and the machine spends its time on \omterm{def:b3:virtual-memory:fault}{page faults} instead of work.
\end{definition}

The second string of \cref{out:b3:virtual-memory:replace} shows the cliff in miniature: a loop over 5 \omterm{def:b3:virtual-memory:paging}{pages} faults on every reference with 4 frames or fewer, under FIFO, LRU and clock alike, and faults only 5 times once 5 frames hold the loop. The remedy for \omterm{def:b3:virtual-memory:working-set}{thrashing} is not a better policy but fewer processes: the kernel suspends some until the \omterm{def:b3:virtual-memory:working-set}{working sets} of the others fit. Peter Denning proposed the \omterm{def:b3:virtual-memory:working-set}{working set} model in 1968. % ledger: b3:denning-1968

\section{Copy on write and memory-mapped files}\label{sec:b3:virtual-memory:cow}

\begin{definition}[Copy on write, memory-mapped file]\label{def:b3:virtual-memory:cow}
\emph{Copy on write}\index{copy on write} shares a \omterm{def:b3:virtual-memory:paging}{page} between two \omterm{def:b3:virtual-memory:space}{address spaces}, mapped read-only in both, and copies it only when one of them writes, in the page-fault handler. A \emph{memory-mapped file}\index{memory-mapped file} is a file whose bytes are mapped into an \omterm{def:b3:virtual-memory:space}{address space}: loads read the file, and stores, in a shared mapping, change it.
\end{definition}

\code{fork} copies the \omterm{def:b3:virtual-memory:space}{address space} of a process; copying every \omterm{def:b3:virtual-memory:paging}{page} would make it slow and, when the child replaces its program at once, useless. Linux implements \code{fork} with \omterm{def:b3:virtual-memory:cow}{copy on write}: what it duplicates is the \omterm{def:b3:virtual-memory:paging}{page tables}. % ledger: b3:linux-fork-cow
\Cref{lst:b3:virtual-memory:cow} measures it with the count of minor \omterm{def:b3:virtual-memory:fault}{page faults}, the faults served without any disk activity, that \code{getrusage} keeps for each process. % ledger: b3:linux-minflt

\omcode{code/bachelor-3/11-virtual-memory/listings/cow.c}{26}{51}{A parent touches 4,096 \omterm{def:b3:virtual-memory:paging}{pages}, forks, and the child reads them, writes them, and writes them again.}\label{lst:b3:virtual-memory:cow}

\omoutput{code/bachelor-3/11-virtual-memory/outputs/cow.txt}{Minor \omterm{def:b3:virtual-memory:fault}{page faults} on the laptop. The first touch of a \omterm{def:b3:virtual-memory:fault}{page faults} once, \omterm{def:b3:virtual-memory:fault}{demand paging}; reading a shared \omterm{def:b3:virtual-memory:paging}{page} costs nothing; the first write to each \omterm{def:b3:virtual-memory:paging}{page} copies it; the second write costs nothing.}\label{out:b3:virtual-memory:cow}

The parent's 4,097 faults are its 4,096 \omterm{def:b3:virtual-memory:paging}{pages}, each mapped on first touch, and one more elsewhere in the program. The child's reads fault zero times: \code{fork} copied the entries, which still point at the parent's frames. Its writes fault once per \omterm{def:b3:virtual-memory:paging}{page}, and each fault copies one \omterm{def:b3:virtual-memory:paging}{page}. The parent still reads 1: the child's writes went to its own copies.

\begin{omfigure}
\begin{tikzpicture}[font=\scriptsize, x=10mm, y=6mm, >=Stealth,
  e/.style={draw, minimum width=15mm, minimum height=5mm, inner sep=1pt},
  f/.style={draw, minimum width=13mm, minimum height=6mm, align=center}]
  \begin{scope}
    \node[font=\small] at (1.6, 2.4) {after \code{fork}};
    \node[e, fill=omDef!10] (p) at (0, 1) {parent: R};
    \node[e, fill=omThm!10] (c) at (0, -0.2) {child: R};
    \node[f, fill=black!6] (fr) at (3, 0.4) {frame 7\\1 1 1 \dots};
    \draw[->] (p.east) -- (fr.west); \draw[->] (c.east) -- (fr.west);
  \end{scope}
  \draw[->, thick] (4.2, 0.4) -- node[above, align=center] {child writes:\\page fault} (6.2, 0.4);
  \begin{scope}[xshift=75mm]
    \node[font=\small] at (1.6, 2.4) {after the copy};
    \node[e, fill=omDef!10] (p2) at (0, 1) {parent: R};
    \node[e, fill=omThm!10] (c2) at (0, -0.2) {child: RW};
    \node[f, fill=black!6] (f7) at (3, 1.2) {frame 7\\1 1 1 \dots};
    \node[f, fill=omThm!12] (f9) at (3, -0.4) {frame 9\\2 1 1 \dots};
    \draw[->] (p2.east) -- (f7.west); \draw[->] (c2.east) -- (f9.west);
  \end{scope}
\end{tikzpicture}
\par\smallskip{\small \omterm{def:b3:virtual-memory:cow}{Copy on write}, for one \omterm{def:b3:virtual-memory:paging}{page}. After \code{fork} both entries point to the same frame, read-only. The child's store faults; the handler copies the frame, points the child's entry at the copy, writable, and the store runs again. The parent's entry stays read-only: its next store faults too, and the handler, finding the frame no longer shared, makes the entry writable without copying.}
\end{omfigure}

A shared mapping of a file turns the file into memory: \code{mmap} returns an address, the \omterm{def:b3:virtual-memory:paging}{pages} are filled from the file on first touch, and a store to a \omterm{def:b3:virtual-memory:paging}{page} marks it dirty; the kernel writes dirty \omterm{def:b3:virtual-memory:paging}{pages} back to the file, at the latest when \code{msync} asks it to. % ledger: b3:linux-map-shared

\omcode{code/bachelor-3/11-virtual-memory/listings/mmap_file.c}{10}{32}{A file edited in place by stores to memory.}\label{lst:b3:virtual-memory:mmap}

\omoutput{code/bachelor-3/11-virtual-memory/outputs/mmap_file.txt}{The file, read back with \code{read}, holds the stores.}\label{out:b3:virtual-memory:mmap}

The same mechanism loads programs: the kernel maps the program file's code \omterm{def:b3:virtual-memory:paging}{pages}, and they come from disk only when executed; and two processes that map the same library share its frames.

\begin{history}[Paging]
L.~A. Belady published a study of replacement algorithms for a virtual-storage computer in 1966. % ledger: b3:belady-1966
Peter Denning published the \omterm{def:b3:virtual-memory:working-set}{working set} model in May 1968. % ledger: b3:denning-1968
Belady, Nelson and Shedler described the anomaly in June 1969, % ledger: b3:belady-nelson-shedler-1969
and Mattson, Gecsei, Slutz and Traiger the \omterm{def:b3:virtual-memory:belady}{stack algorithms} in 1970. % ledger: b3:mattson-1970
\end{history}

\section{Exercises}\label{sec:b3:virtual-memory:exercises}

\begin{exercise}[$\star$]\label{exo:b3:virtual-memory:1}
Split the Sv32 \omterm{def:b3:virtual-memory:space}{virtual address} \code{0x7ffff123}. In the process of \cref{lst:b3:virtual-memory:demo} (root in frame 1, the stack \omterm{def:b3:virtual-memory:paging}{page} in frame 6, its level-0 table in frame 7), which two entry addresses does the walk read, and what is the \omterm{def:b3:virtual-memory:space}{physical address}?
\end{exercise}

\begin{exercise}[$\star$]\label{exo:b3:virtual-memory:2}
A memory access takes 80~ns, a \omterm{def:b3:virtual-memory:tlb}{TLB} lookup no extra time, and a \omterm{def:b3:virtual-memory:tlb}{TLB} miss a walk of two memory reads. What is the average access time with a hit rate of 99~\%, and of 99.9~\%? Which hit rate keeps the slowdown under 5~\%?
\end{exercise}

\begin{exercise}[$\star$]\label{exo:b3:virtual-memory:3}
Count the faults of LRU with 3 frames on 7 0 1 2 0 3 0 4 2 3.
\end{exercise}

\begin{exercise}[$\star$]\label{exo:b3:virtual-memory:4}
\Cref{lst:b3:virtual-memory:demo} adds the line \code{translate(root, 0x10001000, 'r', &pa);}. Predict what it prints.
\end{exercise}

\begin{exercise}[$\star\star$]\label{exo:b3:virtual-memory:5}
A first version of \code{sv32.c} defined its index as below. Loads from the first kilobyte of a mapped \omterm{def:b3:virtual-memory:paging}{page} work; a load from \code{0x10000abc} faults although the \omterm{def:b3:virtual-memory:paging}{page} \code{0x10000000} is mapped. Explain, and fix it.

\omcode{code/bachelor-3/11-virtual-memory/exercises/ex5_vpn_bug.h}{2}{2}{The index as first written.}\label{lst:b3:virtual-memory:vpnbug}
\end{exercise}

\begin{exercise}[$\star\star$]\label{exo:b3:virtual-memory:6}
Write clock with a dirty bit: references are (\omterm{def:b3:virtual-memory:paging}{page}, \code{'r'} or \code{'w'}); to evict, the hand first looks, over one full turn and changing nothing, for a \omterm{def:b3:virtual-memory:paging}{page} neither used nor dirty; failing that, over one more turn, for a \omterm{def:b3:virtual-memory:paging}{page} not used, clearing the bits it passes; and repeats. Count the faults and the write-backs. On (1, w), (2, r), (3, r), (4, r) with 3 frames, which \omterm{def:b3:virtual-memory:paging}{page} does it evict, and which does plain clock evict?
\end{exercise}

\begin{exercise}[$\star\star$]\label{exo:b3:virtual-memory:7}
Prove that OPT makes the fewest possible faults on every reference string, for every number of frames.
\end{exercise}

\begin{exercise}[$\star\star$]\label{exo:b3:virtual-memory:8}
A process has 64~KiB of code at \code{0x00400000}, 8~MiB of heap at \code{0x20000000}, and 64~KiB of stack just below \code{0x80000000}. How many \omterm{def:b3:virtual-memory:paging}{pages} of Sv32 \omterm{def:b3:virtual-memory:paging}{page tables} does it need?
\end{exercise}

\begin{exercise}[$\star\star$]\label{exo:b3:virtual-memory:9}
In \cref{lst:b3:virtual-memory:cow}, the child first writes every other \omterm{def:b3:virtual-memory:paging}{page}, then reads every \omterm{def:b3:virtual-memory:paging}{page}, then writes every \omterm{def:b3:virtual-memory:paging}{page}. Predict the three counts of minor faults.
\end{exercise}

\begin{exercise}[$\star\star\star$]\label{exo:b3:virtual-memory:10}
Prove that LRU is a \omterm{def:b3:virtual-memory:belady}{stack algorithm}.
\end{exercise}

\begin{exercise}[$\star\star\star$]\label{exo:b3:virtual-memory:11}
For any $k$, build a reference string on which clock makes about $k$ times as many faults as LRU with the same number of frames (you choose it).
\end{exercise}

\begin{exercise}[$\star\star\star$]\label{exo:b3:virtual-memory:12}
Sv39 translates 39-bit virtual addresses with 4~KiB \omterm{def:b3:virtual-memory:paging}{pages} through three levels of tables of $2^9$ entries of eight bytes, and also allows 2~MiB megapages and 1~GiB gigapages. % ledger: b3:sv39
Compute the \omterm{def:b3:virtual-memory:tlb}{TLB reach} of 64 entries with each \omterm{def:b3:virtual-memory:paging}{page} size, and the \omterm{def:b3:virtual-memory:paging}{pages} of tables needed to map a 1~GiB array aligned on 1~GiB with small \omterm{def:b3:virtual-memory:paging}{pages}, with megapages, and with one gigapage. Compare with Sv32 in \cref{out:b3:virtual-memory:size}.
\end{exercise}

\section{Weekend problem}\label{sec:b3:virtual-memory:problem}

\begin{problem}[{Weekend problem --- Belady's anomaly}]\label{pb:b3:virtual-memory:1}
More memory should never hurt. This problem checks the claim on FIFO by hand, defines what makes a policy safe, searches every short reference string for counterexamples, and proves that the safe policies are safe.

\omcode{code/bachelor-3/11-virtual-memory/listings/belady_search.py}{5}{37}{Every string up to renaming, extended one reference at a time, with the state of FIFO and LRU with 3 and with 4 frames carried along.}\label{lst:b3:virtual-memory:search}

\omoutput{code/bachelor-3/11-virtual-memory/outputs/belady_search.txt}{The search, lengths 1 to 12, and the anomalous strings found.}\label{out:b3:virtual-memory:search}

\medskip\noindent\textbf{Part I --- FIFO by hand.}
\begin{enumerate}
\item Run FIFO with 3 frames on 1 2 3 4 1 2 5 1 2 3 4 5, writing the frames after each reference. Count the faults.
\item Do the same with 4 frames.
\item From which reference on does the 4-frame run hold \omterm{def:b3:virtual-memory:paging}{pages} the 3-frame run has evicted, and which references then fault only with 4 frames?
\item Clock makes the same counts in \cref{out:b3:virtual-memory:replace}. Why does clock behave exactly as FIFO on this string with 3 frames?
\item Run OPT by hand with 3 frames: which \omterm{def:b3:virtual-memory:paging}{page} does it evict at each fault?
\item What is the fewest faults any policy can make on this string, whatever the number of frames, and why?
\end{enumerate}

\medskip\noindent\textbf{Part II --- The stack property.}
\begin{enumerate}[resume]
\item Give FIFO's \omterm{def:b3:virtual-memory:paging}{pages} with 3 and with 4 frames after 1 2 3 4 1 2 5. Is the inclusion of the definition true?
\item Give LRU's \omterm{def:b3:virtual-memory:paging}{pages} at the same moment.
\item Show that the inclusion at every moment implies that every hit with $f$ frames is a hit with $f + 1$.
\item Can clock be a \omterm{def:b3:virtual-memory:belady}{stack algorithm}? Answer from \cref{out:b3:virtual-memory:replace}.
\item LRU holds the $f$ \omterm{def:b3:virtual-memory:paging}{pages} most recently used. Describe FIFO's \omterm{def:b3:virtual-memory:paging}{pages} in the same way, and say why that description depends on $f$.
\item Explain how a single pass over a string can then give LRU's faults for every number of frames at once.
\end{enumerate}

\medskip\noindent\textbf{Part III --- The search.}
\begin{enumerate}[resume]
\item Why may the search consider only strings in which \omterm{def:b3:virtual-memory:paging}{page} $k + 1$ first appears after \omterm{def:b3:virtual-memory:paging}{page} $k$?
\item The numbers of strings are 1, 2, 5, 15, 52, \dots: check the first four by listing the strings.
\item What is the length of the shortest anomalous string for FIFO with 3 and 4 frames, and how many \omterm{def:b3:virtual-memory:paging}{pages} does it use?
\item How many anomalous strings of that length are there, and what do they share?
\item Some anomalous strings use 6 or 7 \omterm{def:b3:virtual-memory:paging}{pages}. What do the extra \omterm{def:b3:virtual-memory:paging}{pages} do?
\item What does the LRU column show, and why is that no surprise?
\item The count of strings grows faster than any power of a fixed base. Why can the search still afford length 12?
\end{enumerate}

\medskip\noindent\textbf{Part IV --- The proof.}
\begin{enumerate}[resume]
\item Prove \cref{thm:b3:virtual-memory:stack} from question 9.
\item Deduce that LRU never shows the anomaly, using exercise~\ref{exo:b3:virtual-memory:10}.
\item Why does the search not prove that FIFO has no anomaly with 2 and 3 frames?
\item Why does FIFO fail to be a \omterm{def:b3:virtual-memory:belady}{stack algorithm}, in one sentence?
\item Clock shows the anomaly, yet kernels use clock rather than LRU. Why?
\item Sum up: which policies can show the anomaly, how short can an anomalous string be, and why can LRU never show one?
\end{enumerate}
\end{problem}
